\scanpage{91}
\mspage{97}

\uline{Der 3} \textcircled{\scriptsize a} \deleted{\unsure{Soit}}\unsure{Si} $Z\in\mathrm{Diag}$ est $\mathrm{Hot}$-asphérique, alors $Z\xrightarrow{p}e$ est une $\mathbb D$-équivalence, i.e. $\inserted{p^*:}\underset{\shortparallel\, A}{\mathbb D(e)}\unsure{\relbar\joinrel\relbar|\to}\mathbb D(Z)$ est pl. fidèle.

\textcircled{\scriptsize b} (\deleted{\unsure{Renforce}}\inserted{Généralise} \textcircled{\scriptsize a}) Soit $f:Z\to Z'$ une flèche dans $\mathrm{Diag}$. Si $f$ est une $\mathrm{Hot}$-équivalence, c'est une $\mathbb D$-équivalence (i.e. $f^*:\mathbb D(Z')\to\mathbb D(\unsure{Z|2})$ induit une \uline{équivalence} entre les sous-catégories de $Z,Z'$ formées des objets constants)

\textcircled{\scriptsize a'} (Précise considérablement a) Si $Z\in\mathrm{Diag}$ est $\mathrm{Hot}$-asphérique, alors \inserted{on a \textcircled{\scriptsize a}, et} l'image \inserted{essentielle} de $\mathbb D(e)=A$ dans $\mathbb D(Z)$ est formée des objets \uline{localement constants} de $Z$, i.e. \inserted{des objets $F$} tels que le foncteur

\[
\unsure{\relbar}\qquad Z^{\unsure{o|\circ}}\longrightarrow A\qquad z\longmapsto F_z
\]

soit localement constant (i.e. transforme tte flèche en iso), donc \inserted{isom. à un foncteur} constant (puisque $Z$ est $\mathrm{Hot}$-asphérique, donc 1-connexe).

Il faut bien noter que dans a'), il ne \unsure{suffira|suffirait} \uline{pas} de supposer que $Z$ soit $\mathbb D$-asphérique, au lieu de $\mathrm{Hot}$-asphérique. Dès que la $D$-équivalence est \inserted{effectivement} plus faible que $\mathrm{Hot}$, il semble que \textcircled{\scriptsize a'} devienne \inserted{\unsure{toujours}} \deleted{\unsure{faux}} \unsure{f.} si on y \deleted{\unsure{remplace}} suppose seulement que $Z$ soit $\mathbb D$-asphérique. P. ex. même si $\mathbb D=\unsure{\mathbb D_i|\mathbb D_1|D_i}$ (théorie des \unsure{$i$|1}-types d'homotopie \unsure{relative|relatifs}), la \deleted{\unsure{condition}}\inserted{catégorie} \deleted{d'un objet \inserted{local. constant} de $\mathbb D(Z)$}

% slanted
\margin{\unsure{ou sous|ou tout} \unsure{quelles|quel} conditions sur $\mathbb D$ ?}

% archive sheet number: 90
